% =======================================================================
% Relativistic Extension of Chapter X, Section 95
% Effect of Rotation on the Rayleigh--Taylor Instability
%
% Conventions: (-,+,+,+) signature, c explicit, u^mu u_mu = -c^2
% Requires: rel_preamble.tex macros
% =======================================================================

\chapter{Effect of Rotation on the Relativistic Rayleigh--Taylor
  Instability}\label{ch:rel-10-95}

This chapter extends the analysis of \S95 --- the effect of rotation
on the Rayleigh--Taylor (RT) instability --- to the special- and
general-relativistic regime.  When the rotational velocity $\Omega R$
or the gravitational potential become relativistic, the Newtonian
treatment of Chandrasekhar must be replaced by a covariant formulation
in which (i)~the inertial mass density $\rho$ is replaced by the
relativistic enthalpy density $\enthalpy/c^{2} = (\edensity + p)/c^{2}$,
(ii)~the Coriolis force acquires a covariant form tied to the vorticity
of the background 4-velocity, and (iii)~the effective gravity includes
both centrifugal and gravitational contributions, with frame-dragging
corrections in general relativity.  The speed of light~$c$ is kept
explicit so that the classical limit ($c\to\infty$) is recovered
transparently.

% ======================================================================
\section{Relativistic perturbation equations with rotation}
\label{sec:rel-95-eqs}
% ======================================================================

\subsection{Background configuration}

We consider a stratified relativistic perfect fluid in a static
gravitational field, rotating uniformly with angular velocity
$\Omega$ about the vertical ($z$-)~axis.  In the co-rotating frame
the equilibrium 4-velocity is
\begin{equation}
  u^{\mu}_{(0)} = \Lf_{\Omega}\,(c,\,0,\,0,\,0),
  \qquad
  \Lf_{\Omega} = \frac{1}{\sqrt{1 - \Omega^{2}R^{2}/c^{2}}},
  \label{eq:rel95-bg4vel}
\end{equation}
where $R$ is the cylindrical radius.  The fluid is characterised by its
energy density $\edensity(z)$, pressure $p(z)$, and rest-mass density
$\rho_{0}(z)$, with the relativistic enthalpy density
\begin{equation}
  \enthalpy = \edensity + p.
  \label{eq:rel95-enthalpy}
\end{equation}
In the Newtonian limit ($c\to\infty$) we have
$\enthalpy/c^{2}\to\rho_{0}$, and all results below reduce to those
of \S95.

\subsection{Linearised equations in the co-rotating frame}

We denote perturbation velocities by $(u,v,w)$ and the Eulerian
perturbations of pressure and density by $\delta p$ and $\delta\rho$.
In the relativistic co-rotating frame the linearised momentum equations
generalising Chandrasekhar's~\eqref{eq:10-137}--\eqref{eq:10-141}
become
\begin{gather}
  \frac{\enthalpy}{c^{2}}\frac{\partial u}{\partial t}
    - 2\frac{\enthalpy}{c^{2}}\,\Omega\,\Lf_{\Omega}^{2}\,v
    = -\frac{\partial}{\partial x}\,\delta p,
  \label{eq:rel95-mom-x}\\[6pt]
  \frac{\enthalpy}{c^{2}}\frac{\partial v}{\partial t}
    + 2\frac{\enthalpy}{c^{2}}\,\Omega\,\Lf_{\Omega}^{2}\,u
    = -\frac{\partial}{\partial y}\,\delta p,
  \label{eq:rel95-mom-y}\\[6pt]
  \frac{\enthalpy}{c^{2}}\frac{\partial w}{\partial t}
    = -\frac{\partial}{\partial z}\,\delta p
      - g_{\mathrm{eff}}\,\delta\!\left(\frac{\enthalpy}{c^{2}}\right)
      + \sum_{s}T_{s}\!\left(\frac{\partial^{2}}{\partial x^{2}}
        + \frac{\partial^{2}}{\partial y^{2}}\right)\!\delta z_{s}
        \,\delta(z-z_{s}),
  \label{eq:rel95-mom-z}\\[6pt]
  \delta\!\left(\frac{\enthalpy}{c^{2}}\right)
    = -\frac{w}{n}\,D\!\left(\frac{\enthalpy}{c^{2}}\right),
  \label{eq:rel95-cont}\\[6pt]
  \frac{\partial u}{\partial x}
    + \frac{\partial v}{\partial y}
    + \frac{\partial w}{\partial z} = 0.
  \label{eq:rel95-div}
\end{gather}
Three relativistic modifications are visible:

\begin{enumerate}
\item The inertial density $\rho$ is replaced everywhere by
  $\enthalpy/c^{2}$.

\item The Coriolis coupling carries a Lorentz factor
  $\Lf_{\Omega}^{2} = (1-\Omega^{2}R^{2}/c^{2})^{-1}$, which
  enhances the Coriolis force when the co-rotation speed is
  relativistic.  This is the covariant manifestation of the
  relativistic Coriolis effect:
  \begin{equation}
    a^{\mu}_{\text{Cor}}
      = -\frac{2}{c^{2}}\,\epsilon^{\mu\nu\alpha\beta}\,
        u_{\nu}\,\Omega_{\alpha}\,\delta u_{\beta},
    \label{eq:rel95-coriolis}
  \end{equation}
  where $\Omega_{\alpha}$ is the rotation vorticity 4-vector.

\item The density perturbation $\delta\rho$ is replaced by the
  perturbation in $\enthalpy/c^{2}$, and the gravitational
  acceleration is the \emph{effective gravity}~$g_{\mathrm{eff}}$
  (see \S\ref{sec:rel-95-geff} below).
\end{enumerate}

\subsection{Normal-mode reduction}

For perturbations $\propto \exp[i(k_{x}x+k_{y}y)+nt]$ with
$k^{2}=k_{x}^{2}+k_{y}^{2}$, the procedure of \S95 carries through
with the replacements $\rho\to\enthalpy/c^{2}$ and
$\Omega\to\Omega_{\mathrm{rel}}=\Omega\,\Lf_{\Omega}^{2}$.
Define the relativistic $z$-vorticity
\begin{equation}
  \zeta = ik_{y}v - ik_{x}u.
  \label{eq:rel95-vort}
\end{equation}
Combining the $x$- and $y$-equations as in the classical analysis, we
obtain the relativistic counterpart of~\eqref{eq:10-150}:
\begin{equation}
  \zeta = \frac{2\Omega_{\mathrm{rel}}\,Dw}{n},
  \label{eq:rel95-zeta}
\end{equation}
and the counterpart of~\eqref{eq:10-151}:
\begin{equation}
  \frac{\enthalpy}{c^{2}}\,n\!
  \left(1+\frac{4\Omega_{\mathrm{rel}}^{2}}{n^{2}}\right)\!Dw
    = -k^{2}\,\delta p.
  \label{eq:rel95-press}
\end{equation}
Defining the \emph{relativistic effective wavenumber}
\begin{equation}
  \boxed{\;
  \kappa_{\mathrm{rel}}^{2}
    = \frac{k^{2}}{1 + 4\Omega_{\mathrm{rel}}^{2}/n^{2}},
  \qquad
  \Omega_{\mathrm{rel}} = \Omega\,\Lf_{\Omega}^{2}
    = \frac{\Omega}{1-\Omega^{2}R^{2}/c^{2}},\;}
  \label{eq:rel95-kappa}
\end{equation}
the master equation becomes
\begin{equation}
  D\!\left(\frac{\enthalpy}{c^{2}}\,Dw\right)
    - \kappa_{\mathrm{rel}}^{2}\,\frac{\enthalpy}{c^{2}}\,w
    = -\frac{g_{\mathrm{eff}}\,\kappa_{\mathrm{rel}}^{2}}{n^{2}}
      \left\{D\!\left(\frac{\enthalpy}{c^{2}}\right)
        - \frac{k^{2}}{g_{\mathrm{eff}}}\sum_{s}T_{s}\,
          \delta(z-z_{s})\right\}w.
  \label{eq:rel95-master}
\end{equation}
This is formally identical to the non-rotating relativistic RT equation
with $k^{2}$ replaced by $\kappa_{\mathrm{rel}}^{2}$ --- precisely
as in the Newtonian case (\S95, equation~\eqref{eq:10-156}).

% ======================================================================
\section{Relativistic effective gravity: centrifugal, gravitational,
  and frame-dragging contributions}
\label{sec:rel-95-geff}
% ======================================================================

The effective gravitational acceleration felt by the fluid in the
co-rotating frame is, in general relativity,
\begin{equation}
  g_{\mathrm{eff}}
    = g_{\mathrm{grav}}
      + \Omega^{2}R\,\Lf_{\Omega}^{2}
      + g_{\mathrm{fd}},
  \label{eq:rel95-geff}
\end{equation}
where:
\begin{itemize}
\item $g_{\mathrm{grav}}$ is the gravitational acceleration (Schwarzschild:
  $g_{\mathrm{grav}}=GM/(r^{2}\sqrt{1-2GM/rc^{2}})$; in flat spacetime
  we take it as a given external field).

\item $\Omega^{2}R\,\Lf_{\Omega}^{2}$ is the relativistic centrifugal
  acceleration.  In the Newtonian limit this reduces to $\Omega^{2}R$.

\item $g_{\mathrm{fd}}$ is the \emph{frame-dragging} (Lense--Thirring)
  contribution, present when the gravitating body itself rotates.  For
  a Kerr spacetime with angular momentum parameter~$a$,
  \begin{equation}
    g_{\mathrm{fd}} \sim \frac{2GJ}{c^{2}r^{3}}\,\Omega,
    \label{eq:rel95-framedrag}
  \end{equation}
  where $J$ is the angular momentum of the central body.  This term
  has no Newtonian analogue and vanishes as $c\to\infty$.
\end{itemize}

\begin{relcorrection}
The effective gravity~\eqref{eq:rel95-geff} differs from the
Newtonian expression $g + \Omega^{2}R$ in two ways: the centrifugal
term acquires a Lorentz factor $\Lf_{\Omega}^{2}$, diverging as
$\Omega R\to c$; and a frame-dragging term $g_{\mathrm{fd}}$ appears
in general relativity around rotating central masses.  Both
corrections vanish in the limit $c\to\infty$.
\end{relcorrection}

% ======================================================================
\section{Two uniform fluids with rotation: relativistic Coriolis
  modification}
\label{sec:rel-95-twofluids}
% ======================================================================

\subsection{Dispersion relation}

Consider two uniform relativistic fluids of enthalpy densities
$\enthalpy_{1}$ and $\enthalpy_{2}$ ($\enthalpy_{2} > \enthalpy_{1}$
in the lower half-space) separated by a horizontal boundary at
$z=0$, with surface tension~$T$.  Since
equation~\eqref{eq:rel95-master} is formally identical to the
non-rotating case with $k\to\kappa_{\mathrm{rel}}$, we write
down directly the relativistic dispersion relation (cf.\
\S92(a) and equation~\eqref{eq:10-159}):
\begin{equation}
  \boxed{\;
  n^{2} = g_{\mathrm{eff}}\,\kappa_{\mathrm{rel}}\!
  \left[
    \frac{\enthalpy_{2}/c^{2}-\enthalpy_{1}/c^{2}}
         {\enthalpy_{2}/c^{2}+\enthalpy_{1}/c^{2}}
    - \frac{k^{2}T}
           {g_{\mathrm{eff}}\,(\enthalpy_{1}/c^{2}+\enthalpy_{2}/c^{2})}
  \right].
  \;}
  \label{eq:rel95-disp-two}
\end{equation}
Substituting for $\kappa_{\mathrm{rel}}$ from~\eqref{eq:rel95-kappa},
this becomes
\begin{equation}
  n^{2}\!\left(1+\frac{4\Omega_{\mathrm{rel}}^{2}}{n^{2}}\right)^{1/2}
    = g_{\mathrm{eff}}\,k\!
    \left[
      \frac{\enthalpy_{2}-\enthalpy_{1}}
           {\enthalpy_{2}+\enthalpy_{1}}
      - \frac{k^{2}T}
             {g_{\mathrm{eff}}\,(\enthalpy_{1}+\enthalpy_{2})/c^{2}}
    \right].
  \label{eq:rel95-disp-two-exp}
\end{equation}
Define the Atwood number in terms of enthalpy,
\begin{equation}
  \mathcal{A}_{\mathrm{rel}}
    = \frac{\enthalpy_{2} - \enthalpy_{1}}
           {\enthalpy_{2} + \enthalpy_{1}},
  \label{eq:rel95-atwood}
\end{equation}
and let $n_{0,\mathrm{rel}}^{2}$ denote the value of $n^{2}$ in the
absence of rotation (i.e.\ the relativistic RT growth rate without
rotation).  Then equation~\eqref{eq:rel95-disp-two-exp} can be
written as
\begin{equation}
  n^{2}\sqrt{1+4\Omega_{\mathrm{rel}}^{2}/n^{2}}
    = n_{0,\mathrm{rel}}^{2}.
  \label{eq:rel95-n-vs-n0}
\end{equation}
This has exactly the same algebraic structure as the Newtonian
relation~\eqref{eq:10-161}, and its solutions are therefore
\begin{align}
  n^{2} &= -2\Omega_{\mathrm{rel}}^{2}
    + \sqrt{4\Omega_{\mathrm{rel}}^{4}
      + n_{0,\mathrm{rel}}^{4}}
  \quad\text{if } n_{0,\mathrm{rel}}^{2} > 0,
  \label{eq:rel95-sol-unstab}\\[4pt]
  n^{2} &= -2\Omega_{\mathrm{rel}}^{2}
    - \sqrt{4\Omega_{\mathrm{rel}}^{4}
      + n_{0,\mathrm{rel}}^{4}}
  \quad\text{if } n_{0,\mathrm{rel}}^{2} < 0.
  \label{eq:rel95-sol-stab}
\end{align}

\subsection{Stability properties}

From equations~\eqref{eq:rel95-sol-unstab}
and~\eqref{eq:rel95-sol-stab} we draw the same qualitative
conclusions as in \S95:
\begin{enumerate}
\item Rotation does not convert an unstable configuration into a
  stable one, or vice versa: $\operatorname{sgn}(n^{2})
  = \operatorname{sgn}(n_{0,\mathrm{rel}}^{2})$.

\item For stable configurations ($n^{2}<0$), the oscillation
  frequency satisfies
  \begin{equation}
    |n^{2}| \geq 4\Omega_{\mathrm{rel}}^{2}
    \quad\text{as}\quad n_{0,\mathrm{rel}}^{2}\to 0^{-},
    \label{eq:rel95-freq-bound}
  \end{equation}
  i.e.\ the frequency of stable oscillations cannot be less than
  $2\Omega_{\mathrm{rel}}$, the relativistic inertial-wave frequency.

\item The relativistic correction enters through
  $\Omega_{\mathrm{rel}} = \Omega\,\Lf_{\Omega}^{2}>\Omega$.
  Thus relativistic rotation is \emph{more effective} at modifying
  growth rates than Newtonian rotation at the same $\Omega$.
  In the non-relativistic limit ($\Omega R\ll c$),
  $\Omega_{\mathrm{rel}}\to\Omega$ and all results of \S95(a) are
  recovered.
\end{enumerate}

\begin{relcorrection}
Compared with the Newtonian dispersion
relation~\eqref{eq:10-160}, the relativistic version has three
modifications: (1)~$\rho_{1,2}\to\enthalpy_{1,2}/c^{2}$;
(2)~$\Omega\to\Omega_{\mathrm{rel}}=\Omega\,\Lf_{\Omega}^{2}$;
(3)~$g\to g_{\mathrm{eff}}$ including centrifugal Lorentz enhancement
and frame-dragging.  All three vanish in the limit $c\to\infty$.
\end{relcorrection}

% ======================================================================
\section{Exponentially varying density with rotation}
\label{sec:rel-95-exp}
% ======================================================================

\subsection{Relativistic exponential stratification}

We now consider a fluid whose enthalpy density varies exponentially,
\begin{equation}
  \frac{\enthalpy(z)}{c^{2}}
    = \frac{\enthalpy_{0}}{c^{2}}\,e^{\beta z},
  \label{eq:rel95-exp-strat}
\end{equation}
confined between rigid boundaries at $z=0$ and $z=d$.  By the same
reasoning as in \S95(b) (using the formal equivalence
$k\to\kappa_{\mathrm{rel}}$), the relativistic dispersion relation
generalising~\eqref{eq:10-165} is
\begin{equation}
  \frac{g_{\mathrm{eff}}\,\beta}{n^{2}}
    = 1 + \frac{\tfrac{1}{4}\beta^{2}d^{2}+m^{2}\pi^{2}}
               {\kappa_{\mathrm{rel}}^{2}\,d^{2}},
  \label{eq:rel95-exp-disp}
\end{equation}
where $m=1,2,3,\ldots$ is the vertical mode number.  Substituting
$\kappa_{\mathrm{rel}}$ from~\eqref{eq:rel95-kappa},
\begin{equation}
  \frac{g_{\mathrm{eff}}\,\beta}{n^{2}}
    = 1 + \frac{\tfrac{1}{4}\beta^{2}d^{2}+m^{2}\pi^{2}}{k^{2}d^{2}}
      \!\left(1+\frac{4\Omega_{\mathrm{rel}}^{2}}{n^{2}}\right).
  \label{eq:rel95-exp-disp2}
\end{equation}
Solving for $n^{2}$,
\begin{equation}
  \boxed{\;
  n^{2} = n_{0,\mathrm{rel}}^{2}\!
  \left(1 - \frac{4\Omega_{\mathrm{rel}}^{2}}{g_{\mathrm{eff}}\,\beta}
    \cdot\frac{\tfrac{1}{4}\beta^{2}d^{2}+m^{2}\pi^{2}}{k^{2}d^{2}}
  \right),\;}
  \label{eq:rel95-exp-n2}
\end{equation}
where
\begin{equation}
  n_{0,\mathrm{rel}}^{2}
    = \frac{g_{\mathrm{eff}}\,\beta}
           {1+(\tfrac{1}{4}\beta^{2}d^{2}+m^{2}\pi^{2})/(k^{2}d^{2})}
  \label{eq:rel95-n0rel}
\end{equation}
is the non-rotating relativistic growth rate.

\subsection{Modified stability boundaries}

From~\eqref{eq:rel95-exp-n2} we deduce:

\begin{enumerate}
\item If $\beta<0$ (heavy fluid on top in the enthalpy sense), then
  $n_{0,\mathrm{rel}}^{2}<0$ and the stratification is \emph{stable}
  for all wave numbers.  The stable oscillation frequency satisfies
  \begin{equation}
    n^{2}\to n_{0,\mathrm{rel}}^{2}\;\text{for}\;k\to\infty,
    \qquad
    n^{2}\to -4\Omega_{\mathrm{rel}}^{2}\;\text{for}\;k\to 0.
    \label{eq:rel95-stab-limits}
  \end{equation}

\item If $\beta>0$ (unstable stratification), rotation stabilises all
  wave numbers below a critical value~$k_{d,\mathrm{rel}}$ given by
  \begin{equation}
    \boxed{\;
    k_{d,\mathrm{rel}}^{2}\,d^{2}
      = \frac{4\Omega_{\mathrm{rel}}^{2}}
             {g_{\mathrm{eff}}\,\beta}
        \,(\tfrac{1}{4}\beta^{2}d^{2}+m^{2}\pi^{2}).\;}
    \label{eq:rel95-kd}
  \end{equation}
  The minimum critical wave number (over $m$, at $m=1$) is
  \begin{equation}
    k_{d,\mathrm{rel},\min}^{2}
      = \frac{4\Omega_{\mathrm{rel}}^{2}}
             {g_{\mathrm{eff}}\,\beta\,d^{2}}
        \,(\tfrac{1}{4}\beta^{2}d^{2}+\pi^{2}).
    \label{eq:rel95-kd-min}
  \end{equation}
\end{enumerate}

\begin{relcorrection}
Compared with the Newtonian result~\eqref{eq:10-169}, the
relativistic critical wave number~\eqref{eq:rel95-kd} is modified in
two ways:
\begin{itemize}
\item $\Omega^{2}\to\Omega_{\mathrm{rel}}^{2}
  = \Omega^{2}\Lf_{\Omega}^{4}$: relativistic rotation is more
  effective at stabilising large-scale modes.
\item $g\to g_{\mathrm{eff}}$: the effective gravity includes
  centrifugal and frame-dragging terms.
\end{itemize}
As a function of $\beta$, the minimum of $k_{d,\mathrm{rel},\min}$
still occurs at $\beta = 2\pi/d$, exactly as in the Newtonian case.
The relativistic correction factor is
\begin{equation}
  \frac{k_{d,\mathrm{rel}}^{2}}{k_{d}^{2}}
    = \frac{\Omega_{\mathrm{rel}}^{2}}{\Omega^{2}}
      \cdot\frac{g}{g_{\mathrm{eff}}}
    = \frac{\Lf_{\Omega}^{4}\,g}
           {g+\Omega^{2}R\,\Lf_{\Omega}^{2}+g_{\mathrm{fd}}}
    \;\xrightarrow{c\to\infty}\; 1.
  \label{eq:rel95-correction-ratio}
\end{equation}
\end{relcorrection}

% ======================================================================
\section{Modified growth rates}\label{sec:rel-95-growth}
% ======================================================================

For the two-fluid problem, the unstable growth rate
from~\eqref{eq:rel95-sol-unstab} can be written
\begin{equation}
  n_{\mathrm{rel}}^{2}
    = -2\Omega_{\mathrm{rel}}^{2}
      + \sqrt{4\Omega_{\mathrm{rel}}^{4}
        + \bigl(g_{\mathrm{eff}}\,k\,\mathcal{A}_{\mathrm{rel}}
          \bigr)^{2}},
  \label{eq:rel95-growth-explicit}
\end{equation}
where we have set $T=0$ for simplicity.  In the limits:

\begin{itemize}
\item \textbf{No rotation} ($\Omega=0$):
  $n_{\mathrm{rel}}^{2} = g_{\mathrm{eff}}\,k\,
  \mathcal{A}_{\mathrm{rel}}$, the standard relativistic RT growth
  rate.

\item \textbf{Strong rotation}
  ($\Omega_{\mathrm{rel}}^{2}\gg g_{\mathrm{eff}}k\,
  \mathcal{A}_{\mathrm{rel}}$):
  \begin{equation}
    n_{\mathrm{rel}}^{2}
      \approx \frac{(g_{\mathrm{eff}}\,k\,
        \mathcal{A}_{\mathrm{rel}})^{2}}
        {4\Omega_{\mathrm{rel}}^{2}},
    \label{eq:rel95-strong-rot}
  \end{equation}
  so the growth rate is suppressed as $\Omega_{\mathrm{rel}}^{-1}$
  --- stronger suppression than in the Newtonian case where the
  suppression goes as $\Omega^{-1}$.

\item \textbf{Newtonian limit} ($c\to\infty$):
  $\Omega_{\mathrm{rel}}\to\Omega$,
  $g_{\mathrm{eff}}\to g$,
  $\mathcal{A}_{\mathrm{rel}}\to(\rho_{2}-\rho_{1})/(\rho_{2}+\rho_{1})$,
  and~\eqref{eq:rel95-growth-explicit} reduces
  to~\eqref{eq:10-162}.
\end{itemize}

% ======================================================================
\section{Causality check}\label{sec:rel-95-causality}
% ======================================================================

\begin{causalitycheck}
We verify that the perturbation equations and their solutions respect
relativistic causality.

\begin{enumerate}
\item \textbf{Phase velocity bound.}  The phase velocity of
  perturbations is $v_{\mathrm{ph}} = |n|/k$.  From
  equation~\eqref{eq:rel95-growth-explicit}, in the absence of
  surface tension,
  \begin{equation}
    \frac{n^{2}}{k^{2}}
      \leq g_{\mathrm{eff}}\,\frac{\mathcal{A}_{\mathrm{rel}}}{k}
      \leq g_{\mathrm{eff}}\,\frac{1}{k}.
  \end{equation}
  For the RT instability, $n^{2}>0$ corresponds to exponential growth
  (not propagation), so there is no propagating mode whose phase speed
  could exceed~$c$.

\item \textbf{Group velocity bound.}  The stable oscillatory modes
  ($n^{2}<0$, i.e.\ $n=i\omega$) have group velocity
  \begin{equation}
    v_{g} = \frac{\partial\omega}{\partial k}.
  \end{equation}
  From~\eqref{eq:rel95-sol-stab}, in the limit $k\to\infty$,
  $\omega^{2}\to g_{\mathrm{eff}}k$, so
  $v_{g}\sim\sqrt{g_{\mathrm{eff}}/(4k)}\to 0$.  The group velocity
  remains bounded and vanishes for large $k$, well below~$c$.

\item \textbf{Co-rotation speed.}  The formulation requires
  $\Omega R < c$, i.e.\ $\Lf_{\Omega}$ remains finite.  This is
  automatically satisfied for any physical configuration:\ a fluid
  element at cylindrical radius~$R$ cannot co-rotate at or above
  the speed of light.

\item \textbf{Effective wavenumber.}  From~\eqref{eq:rel95-kappa},
  $\kappa_{\mathrm{rel}}^{2}>0$ requires
  $n^{2}>0$ or $n^{2}<-4\Omega_{\mathrm{rel}}^{2}$, which is
  guaranteed by equations~\eqref{eq:rel95-sol-unstab}
  and~\eqref{eq:rel95-sol-stab}.  No modes lie in the forbidden gap
  $-4\Omega_{\mathrm{rel}}^{2}<n^{2}<0$.

\item \textbf{No superluminal signal.}  The system describes an
  incompressible fluid (divergence-free velocity), which is an
  idealisation valid when $v\ll\cs$.  In a fully compressible
  treatment, the fastest signal is the sound speed~$\cs<c$, and the
  RT/rotation corrections do not introduce any new characteristics
  beyond those already present in the relativistic Euler equations.
\end{enumerate}

\textbf{Verdict:} All perturbation modes satisfy $v_{\mathrm{signal}}
\leq c$.  The formulation is causal.
\end{causalitycheck}

% ======================================================================
\section{Summary and non-relativistic limit}
\label{sec:rel-95-summary}
% ======================================================================

The relativistic generalisation of \S95 is achieved by three
systematic replacements:

\begin{center}
\begin{tabular}{lll}
\hline
\textbf{Newtonian} & \textbf{Relativistic} & \textbf{Origin}\\
\hline
$\rho$ & $\enthalpy/c^{2}$ & relativistic inertia\\
$\Omega$ & $\Omega_{\mathrm{rel}}
  = \Omega/(1-\Omega^{2}R^{2}/c^{2})$ & covariant Coriolis\\
$g$ & $g_{\mathrm{eff}}
  = g_{\mathrm{grav}}+\Omega^{2}R\Lf_{\Omega}^{2}+g_{\mathrm{fd}}$
  & centrifugal $+$ frame-dragging\\
\hline
\end{tabular}
\end{center}

\noindent
All dispersion relations, stability criteria, and growth rates of
\S95 carry over with these replacements.  The principal physical
consequence is that relativistic rotation is \emph{more effective}
at modifying the RT instability than its Newtonian counterpart at
the same angular velocity~$\Omega$, owing to the factor
$\Lf_{\Omega}^{2}\geq 1$.

In the non-relativistic limit $c\to\infty$:
$\enthalpy/c^{2}\to\rho$,\;
$\Omega_{\mathrm{rel}}\to\Omega$,\;
$g_{\mathrm{eff}}\to g$,\;
$g_{\mathrm{fd}}\to 0$,
and every equation in this chapter reduces to its counterpart in \S95.
